% Minimal supervisor report for the Purple Team Cybersecurity Lab.
% Export the two Mermaid sources as PNG files using https://mermaid.live,
% place them beside this file, then compile this file with pdfLaTeX.

\documentclass[11pt,a4paper]{article}

\usepackage[T1]{fontenc}
\usepackage[utf8]{inputenc}
\usepackage[margin=2.2cm]{geometry}
\usepackage{graphicx}
\usepackage{booktabs}
\usepackage{tabularx}
\usepackage{enumitem}
\usepackage{xcolor}
\usepackage{hyperref}
\usepackage{microtype}
\usepackage{fancyhdr}

\definecolor{accent}{HTML}{155E75}
\definecolor{light}{HTML}{F1F5F9}

\hypersetup{
  colorlinks=true,
  linkcolor=accent,
  urlcolor=accent,
  pdftitle={Purple Team Cybersecurity Lab -- Supervisor Progress Report},
  pdfauthor={Mariam Ibni}
}

\setlist[itemize]{leftmargin=1.4em,itemsep=2pt,topsep=3pt}
\setlength{\parindent}{0pt}
\setlength{\parskip}{5pt}

\pagestyle{fancy}
\fancyhf{}
\lhead{Purple Team Cybersecurity Lab}
\rhead{Supervisor Progress Report}
\cfoot{\thepage}

\newcommand{\diagram}[2]{%
  \begin{figure}[ht]
    \centering
    \IfFileExists{#1}{%
      \includegraphics[width=0.96\textwidth,height=0.32\textheight,keepaspectratio]{#1}%
    }{%
      \fbox{\parbox[c][3.2cm][c]{0.9\textwidth}{\centering
      Mermaid diagram not rendered yet.\\Run the Mermaid CLI command before compiling.}}%
    }
    \caption{#2}
  \end{figure}
}

\begin{document}

\begin{titlepage}
  \centering
  \vspace*{2.5cm}
  {\Huge\bfseries Purple Team Cybersecurity Lab\par}
  \vspace{0.6cm}
  {\LARGE Supervisor Progress Report\par}
  \vspace{1.5cm}
  \rule{0.75\textwidth}{0.6pt}
  \vspace{1.2cm}

  \begin{tabular}{rl}
    \textbf{Student:} & Mariam Ibni \\
    \textbf{Company:} & CBI \\
    \textbf{Supervisor:} & Yassine Medouar \\
    \textbf{Date:} & 11 August 2026 \\
  \end{tabular}

  \vfill
  \colorbox{light}{%
    \parbox{0.82\textwidth}{%
      \centering
      This short report summarizes the laboratory, the completed offensive and
      defensive work, the current SOAR implementation, and the next priorities.
      Detailed technical evidence remains in the Obsidian project vault.
    }}
  \vfill
\end{titlepage}

\section{Project Objective}

The project builds an isolated Purple Team laboratory in which a controlled
attack is executed across a segmented enterprise-style network and then
detected, investigated, and partially contained using defensive tooling.

The main objectives are to:

\begin{itemize}
  \item demonstrate a realistic attack path from an exposed web application to
        internal Windows and Active Directory systems;
  \item centralize network, web, endpoint, and identity telemetry in Splunk;
  \item develop behavior-based detections mapped to MITRE ATT\&CK;
  \item validate detections through repeatable adversary emulation; and
  \item introduce analyst-approved, temporary, and reversible SOAR response.
\end{itemize}

\section{Laboratory Architecture}

The environment separates the external, DMZ, internal, security-monitoring,
and transit networks. OPNsense provides firewalling and Suricata telemetry,
while Cisco routing connects the internal zones. Splunk centralizes logs and
Shuffle orchestrates response workflows.

\diagram{supervisor-architecture.png}{Simplified laboratory architecture generated from Mermaid.}

\section{Work Completed}

\subsection{Infrastructure}

The following core services are operational:

\begin{itemize}
  \item OPNsense firewall with Suricata and remote syslog;
  \item Cisco inter-zone routing and segmented VLANs;
  \item DVWA in the DMZ;
  \item Active Directory domain controller and domain-joined Windows client;
  \item Splunk with Windows forwarders and OPNsense ingestion; and
  \item self-hosted Shuffle SOAR integrated with Splunk and OPNsense APIs.
\end{itemize}

\subsection{Controlled Offensive Scenario}

The completed attack simulation progressed through:

\begin{enumerate}[leftmargin=1.6em,itemsep=2pt]
  \item external reconnaissance and web discovery;
  \item SQL injection and DVWA command execution;
  \item reverse shell and Chisel tunnelling through the DMZ;
  \item WinRM access to the Windows client;
  \item Active Directory discovery and Kerberoasting;
  \item controlled offline credential recovery and valid-account testing;
  \item authorization-boundary validation;
  \item Run-key persistence and reverse SOCKS simulation; and
  \item a limited password-spray exercise.
\end{enumerate}

All offensive activity was performed in the isolated lab. Secrets and reusable
credential material are excluded from this report and the documentation vault.

\subsection{Detection Engineering}

Splunk detections were built and validated across the full attack path. The
current production catalog contains 21 enabled, application-shared alerts.
Important correlation rules include:

\begin{itemize}
  \item web exploitation followed by an external callback;
  \item long-lived Chisel/WebSocket tunnelling;
  \item DMZ-originated internal discovery and WinRM activity;
  \item successful Windows authentication following remote-management traffic;
  \item the full D13 web-compromise-to-internal-access chain;
  \item Kerberoasting and later service-account use;
  \item Run-key PowerShell persistence and reverse SOCKS execution; and
  \item telemetry-source silence or ingestion delay.
\end{itemize}

D13, D15, and D18 were moved to staggered 15-minute schedules after scheduler
telemetry showed concurrency skips. Post-change validation completed with no
skipped or failed runs and maximum runtimes below 60 seconds.

\diagram{supervisor-lifecycle.png}{Attack, detection, and response lifecycle generated from Mermaid.}

\section{SOAR Playbook PB01}

The first production-candidate playbook is
\texttt{PB01 -- D13 Web Compromise to Internal Access}. Its current flow is:

\begin{itemize}
  \item receive and normalize the D13 result;
  \item enrich Windows authentication and firewall evidence through Splunk;
  \item construct a P1 analyst case;
  \item send an Approve/Reject email through Brevo;
  \item require explicit analyst approval;
  \item validate the alert identity, response mode, entities, and API access;
  \item temporarily add the approved DMZ pivot to an OPNsense external alias;
  \item verify the network effect; and
  \item remove the temporary control and confirm recovery.
\end{itemize}

The containment scope is intentionally narrow: it blocks only new TCP
connections from the approved DMZ pivot to the affected Windows endpoint on
WinRM port 5985. It does not isolate the whole host, disable an account, or
block unrelated services.

\section{Validation Results}

\begin{tabularx}{\textwidth}{@{}p{0.30\textwidth}X@{}}
\toprule
\textbf{Validation} & \textbf{Outcome} \\
\midrule
Splunk to Shuffle & Webhook executions and API enrichment succeeded. \\
Analyst approval & Approve and Reject paths were validated; email delivery succeeded. \\
Dry-run safety gate & Missing live-response mode prevented the firewall write. \\
Live containment & OPNsense added the pivot and a fresh DMZ-to-WinRM connection timed out. \\
Recovery & The alias entry was removed manually and WinRM connectivity returned. \\
Automatic rollback & Not yet accepted; a Shuffle verification branch prevented the delayed rollback from being reached. \\
\bottomrule
\end{tabularx}

The live test therefore proved the containment control but also identified an
important orchestration weakness. A workflow marked \textit{Finished} does not
necessarily mean that every required recovery action completed.

\section{Current Status and Next Steps}

PB01 remains a production candidate. Work is paused before another live test.
The next priorities are:

\begin{enumerate}[leftmargin=1.6em,itemsep=2pt]
  \item schedule rollback directly after a confirmed firewall write so it is
        independent of verification and documentation branches;
  \item test successful, empty-state, API-error, and rollback-error paths;
  \item repeat the complete contain--verify--rollback--recovery cycle;
  \item document and package PB01 as a reusable playbook pattern;
  \item adapt the pattern for reverse SOCKS, Kerberoasting/compromised service
        accounts, and telemetry-loss alerts; and
  \item continue planned endpoint-telemetry and adversary-emulation improvements.
\end{enumerate}

\section{Conclusion}

The project has progressed from infrastructure deployment and adversary
simulation to validated multi-source detections and analyst-controlled response.
The current work demonstrates both the effectiveness of scoped containment and
the importance of testing SOAR failure paths. The remaining PB01 work focuses
on making rollback independent, verifiable, and safe before the design is reused
for additional playbooks.

\end{document}
